%! Function Composition and Decomposition — Unit 1, Lesson 1.3 — AP Practice (ANSWER KEY)
% EXTRA CREDIT — optional, exactly TWO pages, placed between the notes and the graded homework.
% Page 1 is Section I, five multiple-choice items with four options (A)–(D); page 2 is Section
% II, one multi-part free-response question on a pair of tables with interpret-in-context parts.
% Its contexts (a bus's fuel, a wind farm) are used nowhere else in the lesson.
% Distractors are the lesson's errors on purpose: composing in the WRONG ORDER (1A, 2C, 3A, 4B),
% MULTIPLYING f(x)g(x) instead of composing — the 1.2 confusion (1B, 2D, 3C, 5D), and STOPPING
% AT THE HANDOFF and reporting it as the answer (1D, 2A, 5A).
%
% PAGE LOCKSTEP with the other half of the pair: the key marks the correct option in its
% LABEL (no text added to the line) and prose answers are \writespace{H}{…}, fixed height both
% ways. The explicit \newpage fixes the page break in both files.
\documentclass[10pt]{article}
\usepackage{precalculus-article}
\usepackage{precalculus-key}
\newcommand{\TallMath}[1]{$\displaystyle #1\rule[-1.4em]{0pt}{3.2em}$}

\begin{document}

\pageheader{Unit 1, Lesson 1.3}{AP Practice --- Answer Key}

\vspace{0.12in}

\boxguard[8]
\begin{remindbox}
\small
\textbf{Extra credit --- optional.} These two pages are written in the style of the AP
Precalculus exam. Your graded homework is the last section of this packet; do it first. Every
answer choice below is a mistake someone really makes --- one of them stops at the handoff,
one multiplies instead of composing, one runs the machines in the wrong order --- so read all
four before you choose.
\end{remindbox}

\vspace{0.12in}

\begin{headlinebox}{goldbg}
\small
\textbf{Section I --- Multiple Choice.} Circle the letter of the single best answer. No calculator.
\end{headlinebox}

\vspace{0.06in}

\begin{enumerate}[leftmargin=1.9em, itemsep=12pt]

\item The function $g$ gives the number of gallons of fuel a bus uses on a trip of $d$ miles.
      The function $f$ gives the cost, in dollars, of $n$ gallons of fuel. Which expression
      gives the cost of the fuel used on a trip of $d$ miles?
      \begin{enumerate}[leftmargin=2.6em, itemsep=2pt]
        \item[(A)] $g(f(d))$
        \item[(B)] $f(d) \cdot g(d)$
        \item[{\color{keyred}$\checkmark$\,(C)}] $f(g(d))$
        \item[(D)] $g(d)$
      \end{enumerate}

\item \begin{minipage}[t]{0.52\linewidth}
      The table gives every value of the functions $f$ and $g$. What is the value of
      $f(g(3))$?

      \smallskip
      \textbf{(A)} $2$ \hspace{1.6em} {\color{keyred}\textbf{$\checkmark$\,(B)}} $5$ \hspace{1.6em}
      \textbf{(C)} $3$ \hspace{1.6em} \textbf{(D)} $8$
      \end{minipage}\hfill
      \begin{minipage}[t]{0.42\linewidth}
      \vspace{-2pt}
      \centering
      \renewcommand{\arraystretch}{1.3}
      \begin{tabular}{|c|c|c|c|c|}
      \hline
      $x$ & 1 & 2 & 3 & 4 \\
      \hline
      $f(x)$ & 2 & 5 & 4 & 1 \\
      \hline
      $g(x)$ & 4 & 1 & 2 & 3 \\
      \hline
      \end{tabular}
      \end{minipage}

\item Let $f(x) = x^2$ and $g(x) = 2x - 1$. Which of the following is $(f \circ g)(x)$?

      \smallskip
      \textbf{(A)} $2x^2 - 1$ \hspace{1.8em} \textbf{(B)} $4x^2 + 1$ \hspace{1.8em}
      \textbf{(C)} $2x^3 - x^2$ \hspace{1.8em} {\color{keyred}\textbf{$\checkmark$\,(D)}} $4x^2 - 4x + 1$

\item Let $h(x) = \sqrt{3x - 8}$. Which pair of functions satisfies $h(x) = f(g(x))$?
      \begin{enumerate}[leftmargin=2.6em, itemsep=2pt]
        \item[{\color{keyred}$\checkmark$\,(A)}] $f(u) = \sqrt{u}$ \quad and \quad $g(x) = 3x - 8$
        \item[(B)] $f(u) = 3u - 8$ \quad and \quad $g(x) = \sqrt{x}$
        \item[(C)] $f(u) = \sqrt{u} - 8$ \quad and \quad $g(x) = 3x$
        \item[(D)] $f(u) = \sqrt{3u}$ \quad and \quad $g(x) = x - 8$
      \end{enumerate}

\item Let $f(x) = \sqrt{x}$ and $g(x) = x - 9$. Which statement about $f(g(4))$ is true?
      \begin{enumerate}[leftmargin=2.6em, itemsep=2pt]
        \item[(A)] $f(g(4)) = -5$
        \item[{\color{keyred}$\checkmark$\,(B)}] $f(g(4))$ has no value, because the handoff $-5$ is not an input $f$ accepts.
        \item[(C)] $f(g(4)) = -7$, because $\sqrt{4} - 9 = -7$.
        \item[(D)] $f(g(4)) = -10$, because $\sqrt{4} \cdot (4 - 9) = -10$.
      \end{enumerate}

\end{enumerate}

\newpage

\begin{headlinebox}{goldbg}
\small
\textbf{Section II --- Free Response.} Show your reasoning. Answers about the wind farm must be
written \emph{in context}, with units --- that is what earns credit on the AP exam.
\end{headlinebox}

\vspace{0.08in}

\begin{enumerate}[leftmargin=1.9em, itemsep=10pt, start=6]

\item At a wind farm, $S(t)$ is the wind speed in meters per second $t$ hours after midnight,
      and $P(s)$ is the farm's power output in kilowatts when the wind speed is $s$ meters per
      second. Both functions are given entirely by these tables.

      \smallskip
      \renewcommand{\arraystretch}{1.3}
      \begin{center}
      \begin{tabular}{|c|c|c|c|c|c|c|}
      \hline
      $t$ & 0 & 2 & 4 & 6 & 8 & 10 \\
      \hline
      $S(t)$ & 3 & 6 & 9 & 12 & 6 & 18 \\
      \hline
      \end{tabular}
      \qquad
      \begin{tabular}{|c|c|c|c|c|}
      \hline
      $s$ & 3 & 6 & 9 & 12 \\
      \hline
      $P(s)$ & 40 & 150 & 500 & 1100 \\
      \hline
      \end{tabular}
      \end{center}

      \textbf{(a)} Name the handoff in $P(S(4))$, with its units. Then find $P(S(4))$ and
      interpret it in context, with units.
      \writespace{2.0cm}{The handoff is $S(4) = 9$ meters per second. Then $P(9) = 500$: four hours after midnight the farm is producing 500 kilowatts of power.}

      \textbf{(b)} Explain why $S(P(4))$ is not a meaningful quantity at this wind farm.
      \writespace{2.0cm}{$P(4)$ is a power, measured in kilowatts, and $S$ only accepts a time in hours. The handoff would have the wrong units, so $S$ has no question to answer.}

      \textbf{(c)} Find $P(S(2))$ and $P(S(8))$. The two values are equal even though the times
      are different. Explain what that says about the handoffs.
      \writespace{2.2cm}{$S(2) = 6$ and $S(8) = 6$, so both handoffs are 6 meters per second and $P(6) = 150$ kilowatts both times. Different inputs can produce the same handoff, and the same handoff must give the same answer.}

      \textbf{(d)} Can $P(S(10))$ be found from these tables? Answer using the handoff.
      \writespace{2.0cm}{No. The handoff is $S(10) = 18$ meters per second, and the table gives no $P(18)$ --- the outer function does not accept that question, so the composite has no value at $t = 10$.}

      \textbf{(e)} A gearbox is installed so that the blades feel a wind speed 3 meters per
      second higher than the air outside: $b(s) = s + 3$. Write the farm's power output $t$
      hours after midnight as a composition of all three functions, then find its value at
      $t = 2$, naming both handoffs.
      \writespace{2.4cm}{The output is $P(b(S(t)))$. At $t = 2$: the first handoff is $S(2) = 6$ m/s, the second is $b(6) = 9$ m/s, and $P(9) = 500$ kilowatts.}

\end{enumerate}

\end{document}
